% page 98 du fac-simile (image p-097 du scan)
% bloc 160.0 x 237.7 mm ; marge gauche 8.0 mm
\pgc{-12.700mm}{0.000mm}{
\l{\cel{4}90}
\l{}
\l{\sou{ciklo}. I. Periodo kontinua ek ula quanto de yari, dum qua ula}
\l{\cel{3}fenomeni astronomiala eventas sam-ordine. - II. Serio de}
\l{\cel{3}modifikesi di korpo qua pasas tra estadi diversa, e retro-}
\l{\cel{3}venas a sua departo-punto. - III. Serio de poemi pri un}
\l{\cel{3}epoko, un temo, e c. - IV. Lineo spirala quan formacas,}
\l{\cel{3}cirkum stipo o rameto, la origho-punto di folii qui inter-}
\l{\cel{3}korespondas. - DEFIRS.}
\l{}
\l{\sou{ciklameno}. (\sou{bot.}) Planto herbacea, perena, de la familio}
\l{\cel{3}\textquotedbl{}primulacei\textquotedbl{}, di qua la folii esas ronda kordio-forme}
\l{\cel{3}e purpurea, la floro pendanta blanka e purpura. -}
\l{\cel{3}L.\cel{1}\sou{cyclamen}. - DEFIS.}
\l{}
\l{\sou{ciklido}. (\sou{geom.})\cel{2}Areo di la klaso quaresma, qua admisas kom}
\l{\cel{3}lineo duopla la cirklo di la infinito,e posedas dek serii}
\l{\cel{3}de secioni cirkla. - DEFIRS.}
\l{}
\l{\sou{ciklito}. (\sou{patol}.)Varietato di iriso-koroidito qua afektas la}
\l{\cel{3}parto dopa ed extera di la koroido (parto qua havas pro-}
\l{\cel{3}longuri pasable multa). - DEFIS.}
\l{}
\l{\sou{cikloido}.\cel{2}(\sou{geom.}) Kurvo trasata da la par-jiro di punto qua}
\l{\cel{3}apartenas a cirklo-kurvo qua rulas sur rekto fixa. - DEFIRS.}
\l{}
\l{\sou{ciklometrio}\cel{2}Arto mezurar la cirkli. - DEF.}
\l{}
\l{\sou{ciklono}. Sturmo-vento qua vorticas, atraktante (per forco ne-}
\l{\cel{3}rezistebla) irgo quon lu renkontras sur tero e sur maro.}
\l{\cel{3}- DEFIRS.}
\l{}
\l{\sou{ciklopo}. Giganto mitologiala, reprezentata kom havanta nur}
\l{\cel{3}un okulo (meze di sua fronto). - DEFIRS.}
\l{}
\l{\sou{cikloso}. (\sou{biol.})\cel{2}Cirkulado en-celula che la planti (kontraste}
\l{\cel{3}kun la cirkulado generala di la planti).}
\l{}
\l{\sou{cikonio}. (\sou{zool}.). Ucelo migrera, de la klaso \textquotedbl{}steltieri\textquotedbl{}, di}
\l{\cel{3}qua la kolo, la beko, la gambi esas tre longa. - L.cyconia.}
\l{\cel{3}- FISL.}
\l{}
\l{\sou{cikorio}. (\sou{bot.}) Planto legum-gardenala, di qua la folii esas}
\l{\cel{3}mikra e frizita, quan onu manjas kom legumo o salado, de}
\l{\cel{3}qua onu facas infuzuri. - L.\cel{1}\sou{cichorium}. - DEFIRS.}
\l{}
\l{\sou{cikuto}. (\sou{bot.}) Planto perena, venenoza, de la familio \textquotedbl{}umbe-}
\l{\cel{3}liferi\textquotedbl{}. - L.\cel{1}\sou{conium maculatum,}\cel{1}e speco di\cel{1}\sou{cicuta}.) - I.}
\l{}
\l{\sou{cilio}. (\sou{anat.})\cel{2}Pilo qua garnisas la bordo libera di la palpe-}
\l{\cel{3}bro (che la mamiferi). - F. (\sou{bot.}) Sorto di pilo silkatra}
\l{\cel{3}qua bordumas ula parti di la vejetanti. - F.}
\l{}
\l{\sou{ciliaro}. (\sou{anat.}) Parto avana ed extera di la koroido, qua havas}
\l{\cel{3}prolonguri pasable multa.}
}
